\unitopen{05}
\begin{candos}
\item odmienić \de{haben} i powiedzieć, co masz i czego nie masz;
\item użyć biernika (\de{Akkusativ}) w rodzaju męskim: \de{der $\rightarrow$ den}, \de{ein $\rightarrow$ einen};
\item zaprzeczyć rzeczownik (\de{kein}) i czasownik lub przymiotnik (\de{nicht});
\item kupić, zamówić i znaleźć coś: \de{kaufen}, \de{bestellen}, \de{suchen}.
\end{candos}

\section{Słownictwo}
\vocabtable{05}

\section{\de{haben}}
\begin{center}\small
\begin{tabular}{@{}ll@{}}
\de{ich habe} & \de{wir haben}\\
\de{du hast} & \de{ihr habt}\\
\de{er/sie/es hat} & \de{sie/Sie haben}\\
\end{tabular}
\end{center}
\de{haben} łączy się z rzeczownikiem w \emph{bierniku}: \de{Ich habe einen Laptop.} W wyrażeniach
fizycznych polskie „jestem/chce mi się” to niemieckie „mam”: \de{Ich habe Hunger / Durst.}

\section{Biernik: zmienia się tylko rodzaj męski}
Pytanie o dopełnienie: \de{Wen oder was?} (kogo? co?). Po \de{haben}, \de{suchen}, \de{kaufen},
\de{finden}, \de{brauchen}, \de{bestellen}, \de{sehen}, \de{besuchen} rzeczownik stoi w bierniku.

\begin{center}\small
\begin{tabular}{@{}lllll@{}}
\toprule
 & \textsf{\textbf{męski}} & \textsf{\textbf{żeński}} & \textsf{\textbf{nijaki}} & \textsf{\textbf{mnoga}}\\
\midrule
mianownik & \artm{der} / \de{ein} & \artf{die} / \de{eine} & \artn{das} / \de{ein} & \artf{die} / ---\\
biernik & \artm{den} / \de{einen} & \artf{die} / \de{eine} & \artn{das} / \de{ein} & \artf{die} / ---\\
\bottomrule
\end{tabular}
\end{center}

\de{Der Schlüssel ist neu.} $\rightarrow$ \de{Ich suche \textbf{den} Schlüssel.} Rodzaje żeński, nijaki
i liczba mnoga są w bierniku \emph{takie same} jak w mianowniku. Dlatego jedynym nowym wzorcem jest
\artm{den}\,/\,\de{einen}\,/\,\de{keinen}.

\begin{falle}
Polacy znają biernik z polskiego, ale zapominają o nim po niemiecku, bo polska końcówka (\emph{-ę}, \emph{-ą})
nie ma odpowiednika w rzeczowniku: zmienia się \emph{rodzajnik}. Typowy błąd: \de{Ich habe ein Computer.}
Poprawnie: \de{Ich habe \textbf{einen} Computer.} Sprawdzaj przed zdaniem: czy rzeczownik jest męski i czy jest
dopełnieniem?
\end{falle}

\section{Przeczenie: \de{kein} albo \de{nicht}}
\begin{zasada}[Dwa przeczenia, jedna reguła]
\de{kein-} zastępuje \emph{ein-} albo brak rodzajnika przy rzeczowniku:
\de{Ich habe einen Laptop. $\rightarrow$ Ich habe \textbf{keinen} Laptop.} \de{Ich habe Zeit. $\rightarrow$ Ich habe
\textbf{keine} Zeit.} \de{nicht} przeczy resztę: czasownik, przymiotnik, rzeczownik z \emph{der/die/das}
lub z \emph{mein}: \de{Ich kaufe das Auto \textbf{nicht}.} \de{Der Laptop ist \textbf{nicht} neu.}
\end{zasada}
\de{kein-} odmienia się jak \de{ein-}: \de{kein} (męski i nijaki, mianownik), \de{keine}, \de{keinen} (męski, biernik),
\de{keine} (mnoga). \de{nicht} stoi zwykle na końcu zdania lub przed przymiotnikiem/przysłówkiem, który przeczy.

\begin{szwajcaria}
\de{Kein Problem} to w Szwajcarii zwrot kurtuazyjny, często powiedziany zamiast „proszę”. W sklepie
usłyszysz \de{Nei, das han i nöd} (dialekt): to samo co \de{Nein, das habe ich nicht}. Nie musisz mówić
w dialekcie; wystarczy, że to rozpoznasz.
\end{szwajcaria}

\begin{praca}
\de{Ich habe eine Frage zum Code.} \de{Ich habe keinen Zugriff auf den Server.} \de{Das ist kein Bug, das ist ein Feature}
(żart, który zrozumie każdy zespół). Nazwy angielskie zostają, zdanie jest niemieckie.
\end{praca}

\section{Dialogi}
\textbf{Dialog 1 --- w sklepie z elektroniką}
\begin{dialog}
\turn{A}{Guten Tag. Ich suche einen Laptop.}{Dzień dobry. Szukam laptopa.}
\turn{B}{Wir haben einen neuen Laptop. Er ist sehr schnell.}{Mamy nowy laptop. Jest bardzo szybki.}
\turn{A}{Haben Sie auch eine Tastatur?}{Czy mają Państwo też klawiaturę?}
\turn{B}{Ja, wir haben eine Tastatur, aber keine Maus.}{Tak, mamy klawiaturę, ale nie mamy myszy.}
\turn{A}{Kein Problem. Ich kaufe den Laptop und die Tastatur.}{Nie ma problemu. Kupuję laptop i klawiaturę.}
\end{dialog}

\section{Ćwiczenia}
\begin{cwiczenia}
\exitem{Uzupełnij: \de{Du \dots{} einen Computer.}}{\de{hast}}
\exitem{Uzupełnij: \de{Wir \dots{} keine Zeit.}}{\de{haben}}
\exitem{Uzupełnij: \de{Er \dots{} Hunger.}}{\de{hat}}
\exitem{Uzupełnij: \de{Ich suche \dots{} Schlüssel.} (der)}{\de{den}}
\exitem{Uzupełnij: \de{Sie kauft \dots{} Lampe.} (eine)}{\de{eine}}
\exitem{Uzupełnij: \de{Ich brauche \dots{} Stift.} (kein)}{\de{keinen} (\de{der Stift} jest męski)}
\exitem{Uzupełnij: \de{Wir haben \dots{} Problem.} (kein)}{\de{kein} (\de{das Problem} jest nijaki)}
\exitem{Zaprzecz: \de{Ich habe Zeit.}}{\de{Ich habe keine Zeit.}}
\exitem{Zaprzecz: \de{Ich kaufe das Auto.}}{\de{Ich kaufe das Auto nicht.}}
\exitem{Zaprzecz: \de{Der Laptop ist neu.}}{\de{Der Laptop ist nicht neu.}}
\exitem{Popraw: \de{Ich habe ein Computer.}}{\de{Ich habe einen Computer.}}
\exitem{Przetłumacz: „Mam pytanie.”}{\de{Ich habe eine Frage.}}
\exitem{Przetłumacz: „Nie mam hasła.”}{\de{Ich habe kein Passwort.}}
\end{cwiczenia}

\weekend{Opisz po niemiecku pięć rzeczy, które masz, i pięć, których nie masz na biurku
(\de{Ich habe einen Bildschirm. Ich habe keinen Drucker.}). Użyj trzech czasowników z listy, w tym
jednego z biernikiem męskim. Wyślij tekst do korekty.}
